\subsection{Future Work}
\label{sec:roadmap}

The project's phased roadmap continues beyond this document:

\begin{enumerate}
    \item \textbf{Complete the validation suite.} Execute the scripted clock-domain intervention (per-class predictions are stated ex ante: bandwidth classes track the memory clock, compute/on-chip/dependency classes the core clock) and incorporate its verdicts into the artifact.
    \item \textbf{Name the last bytes.} The kernel-argument correlation layer (Layer 3) resolves the two named residual kernels carrying 83\% of unmapped traffic; its target list and join interface are already committed.
    \item \textbf{Embedded-target re-derivation.} Re-run the regime on Jetson-class hardware (the harness already drives remote targets) to re-derive the codegen-dependent quantities above all the spill share and test the taxonomy's transfer to deployment-class hardware.
    \item \textbf{The architecture study.} Feed the traces to Accel-Sim under the generated NDP configurations~\cite{khairy2020accelsim}, with AccelWattch~\cite{kandiah2021accelwattch} for per-class energy, reporting speedup and energy \emph{deltas} against this thesis's measured baselines per the standing rule, never absolutes. The conservative/moderate placement scenarios of \S\ref{sec:syn-model} become simulated designs; the 16--26-kernel offload sets are the experiment list.
    \item \textbf{Breadth via adapters.} Run non-SLAM GPU workloads (DNN inference, GPU database operators) through the unchanged pipeline to test the methodology's generality claim the adapter interface ensures that breadth costs only execution time, not engineering.
\end{enumerate}

The characterization is complete; the design space it opens is measured, bounded, and for the first time for a production Physical-AI perception stack named, byte by byte.
